% Provenance (never printed): NSI terminale "Algorithmique: programmation dynamique (alignement de sequences, rendu
% de monnaie; cout en memoire)". See sources/book2/BRIEFS.md.
\chapter{Dynamic Programming}\label{ch:g12:dynamic-programming}

The recursive \code{fib} of chapter \ref{ch:g12:recursion} needs $2\,692\,537$ calls to compute $\mathrm{fib}(30)$,
a number it could find by hand in $30$ additions. It is not slow because the problem is hard: it is slow because it
solves the same small problems again and again, $\mathrm{fib}(2)$ alone $514\,229$ times, and forgets each answer as
soon as it has used it. Writing the answers down fixes it.

This chapter makes that remedy a method. It applies whenever a problem splits into smaller problems that overlap:
the fewest coins for an amount, where greedy failed; the longest part two texts have in common, and the fewest
changes from one version of a text to the next; the best load of a knapsack with whole objects. Each time the work
is a table filled in the right order, and its price is the memory of the table.

\begin{recall}
A \omterm{def:g12:recursion:recursion}{recursive function} is defined from its values on smaller inputs (chapter \ref{ch:g12:recursion}). A \omterm{def:g11:greedy-algorithms:greedy}{greedy
algorithm} takes the best-looking choice at each step, and fails for some coin systems and for the knapsack with whole
objects (chapter \ref{ch:g11:greedy-algorithms}). A grid is a list of lists, and a dictionary maps keys to values
(Book 1, grades 8 and 9).
\end{recall}

\section{Overlapping subproblems}\label{sec:g12:dynamic-programming:overlap}

\begin{definition}[Subproblem, dynamic programming]\label{def:g12:dynamic-programming:dp}
A \emph{subproblem}\index{subproblem} of a problem is the same problem on a smaller input, whose solution helps build
the solution of the whole. \emph{Dynamic programming}\index{dynamic programming} solves a problem whose subproblems
overlap by solving each subproblem once, recording its solution, and building larger solutions from the recorded
ones.
\end{definition}

\omterm{def:g12:divide-and-conquer:dc}{Divide and conquer} (chapter \ref{ch:g12:divide-and-conquer}) also uses \omterm{def:g12:dynamic-programming:dp}{subproblems}, but the two halves of a \omterm{def:g12:divide-and-conquer:merge}{merge
sort} never share any work. The two calls of \code{fib(n)}, on $n - 1$ and $n - 2$, do: the first one computes
$\mathrm{fib}(n - 2)$ on its way.

\begin{definition}[Memoisation]\label{def:g12:dynamic-programming:memo}
\emph{Memoisation}\index{memoisation} is the recording of a function's results in a table, usually a dictionary from
arguments to results: the function looks the arguments up before computing, and stores the result before returning.
\end{definition}

\omcode{code/grade-12/07-dynamic-programming/listings/fib_memo.py}{1}{35}{The naive recursion and the memoised one,
both counting their calls.}\label{lst:g12:dynamic-programming:fib-memo}

\omoutput{code/grade-12/07-dynamic-programming/outputs/fib_memo.txt}{The memoised version makes $2n - 1$
calls.}\label{out:g12:dynamic-programming:fib-memo}

\begin{omfigure}
\begin{tikzpicture}[font=\scriptsize, box/.style={draw, rounded corners=1pt, inner sep=2pt, fill=omDef!10},
    hit/.style={draw, dashed, rounded corners=1pt, inner sep=2pt, text=black!60}]
  \node[box] (f5) at (0,0) {fib(5)};
  \node[box] (f4) at (-1.2,-0.9) {fib(4)}; \node[hit] (h3) at (1.2,-0.9) {fib(3)};
  \node[box] (f3) at (-2.4,-1.8) {fib(3)}; \node[hit] (h2) at (0,-1.8) {fib(2)};
  \node[box] (f2) at (-3.6,-2.7) {fib(2)}; \node[hit] (h1) at (-1.2,-2.7) {fib(1)};
  \node[box] (f1) at (-4.8,-3.6) {fib(1)}; \node[box] (f0) at (-2.4,-3.6) {fib(0)};
  \foreach \a/\b in {f5/f4, f5/h3, f4/f3, f4/h2, f3/f2, f3/h1, f2/f1, f2/f0} \draw (\a) -- (\b);
  \node[box] at (3.6,-2.4) {computed};
  \node[hit] at (3.6,-3.0) {looked up};
\end{tikzpicture}

{\small The $9$ calls of the memoised \code{fib(5)}, against $15$ for the naive one. The first call on each value
computes it, going down the left side; every later call on it finds it in the dictionary and returns at once, so the
tree has no repeated subtree.}
\end{omfigure}

\begin{omfigure}
\begin{tikzpicture}
\begin{axis}[omrecord, width=11cm, height=5.6cm, ymode=log, xlabel={$n$}, ylabel={calls},
    legend pos=north west, legend style={font=\footnotesize}]
  \addplot[omThm, thick, mark=square*, mark size=1.2pt] table[x=n, y=naive]{code/grade-12/07-dynamic-programming/figures/calls.dat};
  \addplot[omDef, thick, mark=*, mark size=1.2pt] table[x=n, y=memo]{code/grade-12/07-dynamic-programming/figures/calls.dat};
  \legend{naive recursion, memoised}
\end{axis}
\end{tikzpicture}

{\small Calls made by the two versions of \code{fib(n)}, from \code{figures/calls.py}, with a scale that multiplies
by $10$ at each mark. The naive count is multiplied by about $1.6$ at each step and passes $29$ million at $n = 35$;
the memoised count is $2n - 1$, $69$ at $n = 35$.}
\end{omfigure}

Python can keep the notebook itself. The line \code{@cache}, written just above \code{def}, hands the function to
\code{functools.cache}, which returns a version of it that records every result; \code{cache_info} reports how often
the record was used.

\omcode{code/grade-12/07-dynamic-programming/listings/fib_cache.py}{1}{10}{The same \omterm{def:g12:dynamic-programming:memo}{memoisation}, done by the \omterm{def:g11:modules-and-interfaces:library}{standard
library}.}\label{lst:g12:dynamic-programming:fib-cache}

\omoutput{code/grade-12/07-dynamic-programming/outputs/fib_cache.txt}{$91$ values computed, $88$ calls answered from
the record.}\label{out:g12:dynamic-programming:fib-cache}

\section{The bottom-up table}\label{sec:g12:dynamic-programming:bottom-up}

The memoised function still recurses: \code{fib(n)} waits for \code{fib(n - 1)}, which waits for \code{fib(n - 2)},
down to $1$. For large $n$ the chain of waiting calls is longer than Python allows. But the order in which the values
are needed is known in advance: from small to large. Filling the table in that order needs no recursion at all.

\omcode{code/grade-12/07-dynamic-programming/listings/fib_table.py}{1}{24}{Fibonacci bottom-up, with the whole table
and with two values.}\label{lst:g12:dynamic-programming:fib-table}

\omoutput{code/grade-12/07-dynamic-programming/outputs/fib_table.txt}{The table up to $12$; $\mathrm{fib}(90)$ both
ways; $\mathrm{fib}(1000)$, far beyond the recursion limit.}\label{out:g12:dynamic-programming:fib-table}

\begin{method}[Designing a dynamic programme]\label{met:g12:dynamic-programming:design}
(1) Say what one cell of the table means, in words: ``\code{best[a]} is the fewest coins that make \code{a}''.
(2) Write the recurrence: the value of a cell from cells of smaller \omterm{def:g12:dynamic-programming:dp}{subproblems}, and the cells known directly.
(3) Choose an order of filling in which every cell a cell needs is already filled. (4) Read the answer in the last
cell, and, to know how it is achieved and not only its value, follow the choices back. (5) Count the cells: they are
the memory, and, times the work per cell, the time.
\end{method}

Step (5) also shows what can be saved: \code{fib_two} keeps only the last two values, since no cell needs an older
one. The same saving works whenever each row of a table needs only the row before it.

\section{Change with the fewest coins}\label{sec:g12:dynamic-programming:change}

With the coins $1$, $3$ and $4$, greedy pays $6$ as $4 + 1 + 1$, three coins, when $3 + 3$ takes two (chapter
\ref{ch:g11:greedy-algorithms}). The \omterm{def:g12:dynamic-programming:dp}{subproblems} are the smaller amounts. If a best way to pay $a$ uses a coin $c$,
the other coins pay $a - c$, and must do so with as few coins as possible: otherwise a better way to pay $a - c$
would give a better way to pay $a$. So
\[
\mathrm{best}(0) = 0, \qquad \mathrm{best}(a) = 1 + \min_{c \le a} \mathrm{best}(a - c) \quad (a \ge 1),
\]
the minimum running over the coins $c$ of the system that are at most $a$.

\omcode{code/grade-12/07-dynamic-programming/listings/change_dp.py}{1}{32}{The table of fewest coins, and the coins
rebuilt from the recorded choices.}\label{lst:g12:dynamic-programming:change}

\omoutput{code/grade-12/07-dynamic-programming/outputs/change_dp.txt}{Two coins for $6$ where greedy used three; and
$36$ in Brenmoor, the smallest amount greedy got wrong there.}\label{out:g12:dynamic-programming:change}

\begin{omfigure}
\begin{tikzpicture}[font=\footnotesize, x=11mm, y=6mm]
  \foreach \a/\b/\l in {0/0/{--}, 1/1/1, 2/2/1, 3/1/3, 4/1/4, 5/2/1, 6/2/3} {
    \node[draw, minimum width=11mm, minimum height=6mm, fill=black!5] (a\a) at (\a,1) {$\a$};
    \node[draw, minimum width=11mm, minimum height=6mm] (b\a) at (\a,0) {$\b$};
    \node[draw, minimum width=11mm, minimum height=6mm] at (\a,-1) {$\l$};
  }
  \node[left] at (-0.6,1) {amount}; \node[left] at (-0.6,0) {\code{best}}; \node[left] at (-0.6,-1) {\code{last}};
  \foreach \a in {0, 3, 6} \draw[omThm, very thick] (\a-0.5,-1.5) rectangle (\a+0.5,1.5);
  \draw[->, thick, omThm] ([yshift=4mm]a6.north) to[bend right=35] node[above, font=\scriptsize] {coin $3$}
    ([yshift=4mm]a3.north);
  \draw[->, thick, omThm] ([yshift=4mm]a3.north) to[bend right=35] node[above, font=\scriptsize] {coin $3$}
    ([yshift=4mm]a0.north);
\end{tikzpicture}

{\small The table for the coins $1$, $3$, $4$, filled from left to right. \code{best[6]} $= 2$ came from
\code{best[3]} $+ 1$, and \code{last[6]} records the coin $3$; the rebuild follows the recorded coins back from $6$ to
$3$ to $0$.}
\end{omfigure}

The table has $a + 1$ cells, and each looks at every coin: for $k$ coins, about $a \times k$ steps. Trying all
combinations of coins, as in chapter \ref{ch:g11:greedy-algorithms}, took $3687$ combinations for $100$ in Brenmoor;
the table takes about $600$ steps.

\section{Common subsequences and alignments}\label{sec:g12:dynamic-programming:lcs}

\begin{definition}[Longest common subsequence]\label{def:g12:dynamic-programming:lcs}
A \emph{subsequence} of a sequence keeps some of its elements, in their order, and drops the others: \code{LIBRY} is a
subsequence of \code{LIBRARY}. A \emph{longest common subsequence}\index{longest common subsequence} of two sequences
is a sequence that is a subsequence of both, and as long as possible.
\end{definition}

Let $t[i][j]$ be the length of a \omterm{def:g12:dynamic-programming:lcs}{longest common subsequence} of the first $i$ letters of $a$ and the first $j$ of $b$.
If $a$'s letter $i$ and $b$'s letter $j$ are equal, one such subsequence ends with that letter, and
$t[i][j] = t[i - 1][j - 1] + 1$. If they differ, at least one of them is not the last letter of the subsequence and can
be dropped: $t[i][j] = \max(t[i - 1][j], t[i][j - 1])$. A row or a column with no letters holds $0$. Filling row by row
from the top left, each cell finds the three cells it needs above it and on its left.

\omcode{code/grade-12/07-dynamic-programming/listings/lcs.py}{1}{35}{The table of common subsequences, and a longest
one rebuilt.}\label{lst:g12:dynamic-programming:lcs}

\omoutput{code/grade-12/07-dynamic-programming/outputs/lcs.txt}{\code{LIBRARY} and \code{LITERARY} share
\code{LIRARY}, six letters.}\label{out:g12:dynamic-programming:lcs}

\begin{omfigure}
\begin{tikzpicture}[font=\footnotesize, x=6.5mm, y=-6.5mm]
  \foreach \c [count=\j] in {L,I,T,E,R,A,R,Y} \node at (\j,-1) {\c};
  \foreach \c [count=\i] in {L,I,B,R,A,R,Y} \node at (-1,\i) {\c};
  \foreach \i/\j in {0/0, 1/1, 2/2, 2/3, 2/4, 3/4, 4/5, 5/6, 6/7, 7/8}
    \fill[omMeth!18] (\j-0.5,\i-0.5) rectangle (\j+0.5,\i+0.5);
  \foreach \row [count=\i from 0] in {{0,0,0,0,0,0,0,0,0}, {0,1,1,1,1,1,1,1,1}, {0,1,2,2,2,2,2,2,2},
      {0,1,2,2,2,2,2,2,2}, {0,1,2,2,2,3,3,3,3}, {0,1,2,2,2,3,4,4,4}, {0,1,2,2,2,3,4,5,5},
      {0,1,2,2,2,3,4,5,6}}
    \foreach \v [count=\j from 0] in \row \node at (\j,\i) {$\v$};
  \foreach \x in {-0.5,0.5,...,8.5} \draw[black!30] (\x,-0.5) -- (\x,7.5);
  \foreach \y in {-0.5,0.5,...,7.5} \draw[black!30] (-0.5,\y) -- (8.5,\y);
  \foreach \i/\j in {1/1, 2/2, 4/5, 5/6, 6/7, 7/8} \draw[omThm, thick] (\j,\i) circle (2.6mm);
\end{tikzpicture}

{\small The table of \code{LIBRARY} (rows) and \code{LITERARY} (columns). The shaded cells are the path of the
rebuild, from the bottom right back to the top left; at each circled cell the two letters are equal, and that letter
belongs to the subsequence: L, I, R, A, R, Y.}
\end{omfigure}

A second question about two versions of a text is how many changes turn one into the other.

\begin{definition}[Edit distance, alignment]\label{def:g12:dynamic-programming:edit-distance}
The \emph{edit distance}\index{edit distance} between two sequences is the smallest number of operations that turn the
first into the second, each operation inserting one element, deleting one, or replacing one by another. An
\emph{alignment}\index{alignment} of two sequences writes them one above the other, with gaps, so that each column
holds two equal elements, two different ones (a replacement), or an element and a gap (an insertion or a deletion).
\end{definition}

The table has the same shape. $d[i][j]$ is the distance from the first $i$ elements of $a$ to the first $j$ of $b$. In
a best \omterm{def:g12:dynamic-programming:edit-distance}{alignment}, the last column either pairs $a$'s element $i$ with $b$'s element $j$, for no cost if they are equal
and $1$ otherwise; or holds $a$'s element $i$ over a gap, a deletion; or a gap over $b$'s element $j$, an insertion.
So $d[i][j]$ is the smallest of $d[i - 1][j - 1]$ (plus $1$ if the elements differ), $d[i - 1][j] + 1$ and
$d[i][j - 1] + 1$; and $d[i][0] = i$, $d[0][j] = j$. The elements can be letters, or the words of a sentence.

\omcode{code/grade-12/07-dynamic-programming/listings/align.py}{1}{37}{The distance table, and an \omterm{def:g12:dynamic-programming:edit-distance}{alignment}
rebuilt from it.}\label{lst:g12:dynamic-programming:align}

\omcode{code/grade-12/07-dynamic-programming/listings/align.py}{39}{56}{Printing an \omterm{def:g12:dynamic-programming:edit-distance}{alignment} of two
sentences, word by word.}\label{lst:g12:dynamic-programming:align-show}

\omoutput{code/grade-12/07-dynamic-programming/outputs/align.txt}{Three changes: one word inserted, two words
replaced. The marks are \code{|} for a word kept and \code{*} for a word replaced.}\label{out:g12:dynamic-programming:align}

Several \omterm{def:g12:dynamic-programming:edit-distance}{alignments} can reach the same distance. The rebuild keeps an element whenever it can, and otherwise prefers a
gap to a replacement, which keeps the \omterm{def:g12:dynamic-programming:edit-distance}{alignment} close to what a reader would write.

\section{The knapsack with whole objects}\label{sec:g12:dynamic-programming:knapsack}

A knapsack holds at most $W$ kilograms; each object has a mass and a value, and is taken whole or left. Greedy by
value per kilogram fails here (chapter \ref{ch:g11:greedy-algorithms}): with $10$ kilograms, it takes the object of
$6$~kg worth $30$ and then nothing fits, when two objects of $5$~kg worth $20$ each give $40$. Trying every subset is
exact, but there are $2^n$ of them.

The \omterm{def:g12:dynamic-programming:dp}{subproblems} use fewer objects and less room. Let $k[i][w]$ be the best value with the first $i$ objects and at
most $w$ kilograms. Object $i$ is left, and then $k[i][w] = k[i - 1][w]$; or, if it fits, taken, and then the rest is
the best load of $w$ minus its mass with the objects before it. The table is filled row by row, each row from the one
above.

\omcode{code/grade-12/07-dynamic-programming/listings/knapsack_dp.py}{1}{31}{The knapsack table, the objects rebuilt,
and every subset as a check.}\label{lst:g12:dynamic-programming:knapsack}

\omcode{code/grade-12/07-dynamic-programming/listings/knapsack_dp.py}{32}{42}{The small case of the greedy chapter,
then a pack of eight objects.}\label{lst:g12:dynamic-programming:knapsack-run}

\omoutput{code/grade-12/07-dynamic-programming/outputs/knapsack_dp.txt}{The table of the small case, row $i$ for the
first $i$ objects; then $240$ in $15$~kg, as every subset finds.}\label{out:g12:dynamic-programming:knapsack}

\begin{omfigure}
\begin{tikzpicture}[font=\footnotesize, x=7mm, y=-6mm]
  \foreach \w in {0,...,10} \node[black!60] at (\w,-1) {$\w$};
  \foreach \t [count=\i from 0] in {--, A\,6\,kg\,30, B\,5\,kg\,20, C\,5\,kg\,20} \node[left] at (-0.6,\i) {\t};
  \fill[omMeth!18] (9.5,2.5) rectangle (10.5,3.5); \fill[omMeth!18] (4.5,1.5) rectangle (5.5,2.5);
  \fill[omMeth!18] (-0.5,0.5) rectangle (0.5,1.5);
  \foreach \row [count=\i from 0] in {{0,0,0,0,0,0,0,0,0,0,0}, {0,0,0,0,0,0,30,30,30,30,30},
      {0,0,0,0,0,20,30,30,30,30,30}, {0,0,0,0,0,20,30,30,30,30,40}}
    \foreach \v [count=\w from 0] in \row \node at (\w,\i) {$\v$};
  \foreach \x in {-0.5,0.5,...,10.5} \draw[black!30] (\x,-0.5) -- (\x,3.5);
  \foreach \y in {-0.5,0.5,...,3.5} \draw[black!30] (-0.5,\y) -- (10.5,\y);
  \draw[->, thick, omThm] (9.6,2.9) -- (5.4,2.1);
  \draw[->, thick, omThm] (4.6,1.9) -- (0.4,1.1);
  \foreach \t [count=\i] in {A left, B taken, C taken} \node[font=\scriptsize, right, omThm] at (10.7,\i) {\t};
  \node[font=\scriptsize, black!60, left] at (-0.6,-1) {mass $w$:};
\end{tikzpicture}

{\small The table of the small case, objects in rows, mass allowed in columns. $40$ in the last cell differs from the
$30$ above it, so C was taken, leaving $5$~kg; $20$ in row B differs from the $0$ above it, so B was taken, leaving
$0$~kg; A was left.}
\end{omfigure}

The table has $(n + 1)(W + 1)$ cells and each takes a few steps: for $8$ objects and $15$~kg, $144$ cells against
$256$ subsets. The cost grows with $W$, the size of the knapsack, not only with the number of objects: masses to the
gram in a $15$~kg knapsack make $15\,001$ columns. Tables of two whole texts grow faster still.

\omcode{code/grade-12/07-dynamic-programming/listings/table_sizes.py}{1}{9}{The memory of an $n \times n$ table, and
of two rows.}\label{lst:g12:dynamic-programming:sizes}

\omoutput{code/grade-12/07-dynamic-programming/outputs/table_sizes.txt}{Two texts of a million words would need eight
terabytes for the whole table, and $16$ megabytes for two rows.}\label{out:g12:dynamic-programming:sizes}

The distance needs only the row above: two rows are enough, as for \code{fib_two}. The \omterm{def:g12:dynamic-programming:edit-distance}{alignment} does not: the rebuild
walks back through the whole table. Keeping both the small memory and the \omterm{def:g12:dynamic-programming:edit-distance}{alignment} is possible, by a cleverer
method that is out of reach here.

\begin{history}[Bellman, 1954]
Richard Bellman published ``The theory of \omterm{def:g12:dynamic-programming:dp}{dynamic programming}'', an early general account of the method, in 1954 in
the \emph{Bulletin of the American Mathematical Society}. In its title, \emph{programming} means planning a sequence
of decisions, as in a schedule, not writing code.
% ledger: g12:bellman-1954
\end{history}

\section{Exercises}\label{sec:g12:dynamic-programming:exercises}

\begin{exercise}[$\star$]\label{exo:g12:dynamic-programming:1}
How many calls does \code{fib_memo(6)} make? And \code{fib_naive(6)}?
\end{exercise}

\begin{exercise}[$\star$]\label{exo:g12:dynamic-programming:2}
Fill the table \code{best} for the coins $1$, $3$, $4$ up to the amount $7$. Which coins pay $7$?
\end{exercise}

\begin{exercise}[$\star$]\label{exo:g12:dynamic-programming:3}
Find a \omterm{def:g12:dynamic-programming:lcs}{longest common subsequence} of \code{STATION} and \code{NATIONS}, and its length.
\end{exercise}

\begin{exercise}[$\star$]\label{exo:g12:dynamic-programming:4}
In the knapsack table of the small case, from which cells is $k[2][10] = 30$ computed?
\end{exercise}

\begin{exercise}[$\star$]\label{exo:g12:dynamic-programming:5}
Write \code{binomial(n, k)}, memoised, from Pascal's rule $\binom{n}{k} = \binom{n-1}{k-1} + \binom{n-1}{k}$, with
$\binom{n}{0} = \binom{n}{n} = 1$.
\end{exercise}

\begin{exercise}[$\star\star$]\label{exo:g12:dynamic-programming:6}
A student memoises \code{binomial} with \code{memo[n]} instead of \code{memo[(n, k)]}. The call
\code{binomial_bad(10, 5)} returns $96$. Explain.
\end{exercise}

\begin{exercise}[$\star\star$]\label{exo:g12:dynamic-programming:7}
Write \code{min_coins(amount, coins)} top-down, as a memoised recursion, returning \code{None} when the amount cannot
be paid.
\end{exercise}

\begin{exercise}[$\star\star$]\label{exo:g12:dynamic-programming:8}
A student fills the table of \code{lcs_table} with \code{range(n, 0, -1)} and \code{range(m, 0, -1)}, the same
recurrence otherwise. On \code{LIBRARY} and \code{LITERARY} the answer is $1$. Why?
\end{exercise}

\begin{exercise}[$\star\star$]\label{exo:g12:dynamic-programming:9}
How many paths go from the top-left corner to the bottom-right corner of a grid of $5 \times 5$ squares, along its
lines, moving only right or down? Set up a table.
\end{exercise}

\begin{exercise}[$\star\star$]\label{exo:g12:dynamic-programming:10}
Write \code{diff_lines(old, new)}, which lists the lines of two versions of a file: lines kept with two spaces in
front, lines deleted with \code{-}, lines added with \code{+}, using a \omterm{def:g12:dynamic-programming:lcs}{longest common subsequence} of the lines.
\end{exercise}

\begin{exercise}[$\star\star$]\label{exo:g12:dynamic-programming:11}
The program of the chapter prints the lines of \code{show}, joined by spaces. What does it print for the letters of
\code{kitten} and \code{sitting}? What is their distance?
\end{exercise}

\begin{exercise}[$\star\star\star$]\label{exo:g12:dynamic-programming:12}
Prove, by induction on $a$, that \code{change_table} gives the fewest coins for every amount $a$ that can be paid.
\end{exercise}

\begin{exercise}[$\star\star\star$]\label{exo:g12:dynamic-programming:13}
Write \code{edit_operations(a, b)}, the list of operations of a best \omterm{def:g12:dynamic-programming:edit-distance}{alignment}, such as \code{"replace k by s"}.
Check that it has as many operations as the distance.
\end{exercise}

\begin{exercise}[$\star\star\star$]\label{exo:g12:dynamic-programming:14}
Count the cells and the steps of \code{knapsack} for $n$ objects and a capacity $W$. A removal van holds $3$ tonnes;
its $40$ boxes are weighed to the kilogram, then to the gram. How many cells in each case? What does it say about the
cost?
\end{exercise}

\begin{exercise}[$\star\star\star$]\label{exo:g12:dynamic-programming:15}
Write \code{longest_increasing(t)}, which returns a longest strictly increasing subsequence of a list, with a table
of the longest one ending at each position and a rebuild. Count its steps.
\end{exercise}

\section{Problem: Aligning Two Versions}\label{sec:g12:dynamic-programming:problem}

\begin{problem}[{Weekend problem --- what changed in the notice?}]\label{pb:g12:dynamic-programming:1}
The library on Town Square has rewritten its notice. The program compares the two versions word by word, with the
tables of the chapter, and prints only the columns of the \omterm{def:g12:dynamic-programming:edit-distance}{alignment} that changed.

\omcode{code/grade-12/07-dynamic-programming/listings/versions.py}{1}{37}{The two versions, compared word by
word.}\label{lst:g12:dynamic-programming:versions}

\omoutput{code/grade-12/07-dynamic-programming/outputs/versions.txt}{Words, common words, distance, the changed
columns, and the size of the tables.}\label{out:g12:dynamic-programming:versions}

\needspace{4\baselineskip}\medskip\noindent\textbf{Part I --- The two versions.}
\begin{enumerate}
  \item Read the two notices. List, in your own words, what the library changed.
  \item Why does the program compare words and not letters?
  \item \code{Sunday.} and \code{Sunday} count as different words. Why? What would change if the program removed the
    punctuation first?
\end{enumerate}

\needspace{4\baselineskip}\medskip\noindent\textbf{Part II --- Common words.}
\begin{enumerate}[resume]
  \item What does \code{lcs_table(a, b)[n][m]} measure here?
  \item How many words of the old version are not in the common subsequence? Of the new one?
  \item Why is the \omterm{def:g12:dynamic-programming:edit-distance}{edit distance} at least the larger of these two numbers, and at most their sum?
\end{enumerate}

\needspace{4\baselineskip}\medskip\noindent\textbf{Part III --- The distance and the alignment.}
\begin{enumerate}[resume]
  \item Count the replacements, deletions and insertions in the printed columns. Check the distance.
  \item The deleted words \code{cents a day.} come after \code{ten -> nothing.}. Give another \omterm{def:g12:dynamic-programming:edit-distance}{alignment} of the same
    cost for these words.
  \item Why is \code{appointment -> appointment.} a replacement?
  \item \code{Thursday} and \code{and} are deleted. Write the columns of the \omterm{def:g12:dynamic-programming:edit-distance}{alignment} around them.
  \item Each row of the distance table is filled from left to right. Which cells does $d[i][j]$ need, and why is this
    order right?
  \item The rebuild starts at the bottom-right cell. Why does it never leave the table, and why does it stop?
  \item Prove that a column holding two different words in the rebuilt \omterm{def:g12:dynamic-programming:edit-distance}{alignment} costs one change, and that the
    total of the costs is $d[n][m]$.
\end{enumerate}

\needspace{4\baselineskip}\medskip\noindent\textbf{Part IV --- Memory.}
\begin{enumerate}[resume]
  \item How many cells has the full table? How many bytes, at $8$ bytes per cell?
  \item Explain \code{distance_two_rows}: what does \code{previous} hold at the start of each pass of the loop?
  \item Why can the two-row version give the distance but not the \omterm{def:g12:dynamic-programming:edit-distance}{alignment}?
  \item Two versions of a book of $100\,000$ words each: memory of the full table, and of two rows?
  \item Time: how many cells are computed in each version? Did saving memory cost time?
  \item The tests compare \code{distance_two_rows} with the full table on $150$ random pairs. Why is that a
    convincing check?
  \item State the result: the distance between the two notices, the common words, and the cells of the two tables.
\end{enumerate}
\end{problem}
